\subsubsection*{命题 1.13.}
\underline{若 $X_s$ 在 $x_0$ 处有普通二次奇点, 则猜想 1.9 对 $(X/S,x_0)$ 成立.}

\medskip
\noindent\underline{情形 1: $n=\dim(X_s)=2m-1$.}

此问题在 $X$ 的平展拓扑下是局部的. 因此可假设 $X$ 为 $\mathbb E_S^{n+1}$ 中由方程
\[
f=\sum_{i=1}^{m}x_i x_{i+m}+a=0,
\]
定义的闭子概形, 其中 $a$ 属于 $V$ 的极大理想, $x_0$ 为坐标原点. 此外, $X$ 在 $x_0$ 处正则当且仅当 $a$ 为一致化参数. 由 1.4,
\[
\mu=\operatorname{length}\!\left(
\mathcal O_{\mathbb E_S^{2m},x_0}/
\left(f,\frac{\partial f}{\partial x_i}\right)\right)=1.
\]
由第 XV 讲对普通二次奇点的消灭圈的计算, 当 $i\ne n$ 时 $\Phi^i_{x_0}=0$, 而 $\Phi^n_{x_0}$ 的维数为一, 且惯性群在其上平凡作用.\footnote{原文在所印的“由”字之后, 将应有的引文留空. 此处所述的集中性和惯性作用, 正是第 XV 讲已建立的普通二次奇点的消灭圈计算.}
因此消灭圈的野数为零, 总数为 $1$, 这验证了 1.9.

\medskip
\noindent\underline{情形 2: $n=\dim(X_s)=2m$ 且 $\operatorname{char}(k(s))\ne2$.}

在这些假设下, 可取 $X\subset\mathbb E_S^{n+1}$ 由
\[
f=\sum_{i=1}^{m}x_i x_{i+m}+x_{2m+1}^2+a=0,
\]
定义, 其中 $a$ 为一致化参数. 仍有 $\mu=1$, 当 $i\ne n$ 时 $\Phi^i_{x_0}=0$, 且 $\Phi^n_{x_0}$ 的维数为一. 此时惯性群通过 $\pm1$ 作用于 $\Phi^n_{x_0}$, 因而作用温和. 消灭圈的野数仍为零, 从而验证了 1.9.

\medskip
\noindent\underline{情形 3: $n=\dim(X_s)=2m$ 且 $\operatorname{char}(k(s))=2$.}

此时可取 $X\subset\mathbb E_S^{n+1}$ 由
\[
f=\sum_{i=1}^{m}x_i x_{i+m}
 +(x_{2m+1}^2+a x_{2m+1}+b)=0,
\]
定义, 其中 $a$ 和 $b$ 均属于极大理想. 概形 $X$ 正则当且仅当 $b$ 为一致化参数, 其泛纤维光滑当且仅当 $a^2-4b\ne0$.

$X/S$ 的 Milnor 数与由
\begin{equation}\tag{1}
x^2+a x+b=0.
\end{equation}
定义的 $X_0\subset\mathbb E_S^1$ 的 Milnor 数相同. 同样, 当 $i\ne n$ 时 $\Phi^i_{x_0}=0$, $\Phi^n_{x_0}$ 的维数为一, 且 $I$ 通过 $\pm1$ 作用于其上; 作用平凡的指数为二的 $I$ 的子群, 对应于 (1) 所定义的 $K$ 的扩张. 因此 $X/S$ 在 $x_0$ 处的消灭圈总数等于 $X_0/S$ 的消灭圈总数, 而此情形的 1.9 由 1.12 推出. Milnor 数由 1.5 的判别式公式给出:
\begin{equation}\tag{1.13.1}
\mu=v(a^2-4b).
\end{equation}
